\sscp{Retention and loss of final consonants}{07-02-01}
\citet{Stresemann1927}, \citet{Collins1981,Collins1983}
and \citet{Blust1981} note that while many historical word-final
consonants are lost in central Maluku,
a few are retained.
\citet{Grimes2024} notes that PMP final consonants *n,
*t, *R are retained on some nouns (often relating to nature).
\trf{07-01} shows retentions of final *n across several subgroups.
In these tables
and all subsequent tables in this chapter dividing lines separate
the subgroups we identify clockwise from west to \ili{east} to south:
\tsc{Taliabu}, \tsc{Sula-Buru}, \tsc{\ili{Ambon}-Seram}, \ili{Banda}, and \tsc{Seram-Tanimbar-Bomberai}.
Some tables do not have data from each of these five subgroups.
%XXX put this info about tables somewhere in the intro to chapter? Or do you \it{really}
%want the text spanning columns explicitly marking the subgroups?
%I removed it because in many cases we had to split tables into
%2 and/or are really squeezing for space.}

\begin{table}[h]\tcp{Geographical central Maluku retention of *n /{\gap\#}}{07-01}
	\begin{threeparttable}\st{6pt}
	\begin{tabular}{lllll}\lsptoprule
*gloss	&	wind	&	moon	&	rain	&	path, road\\
PMP	&	*haŋi\tbr{n}	&	*bula\tbr{n}	&	*quza\tbr{n}	&	*zala\tbr{n}\\
%\midrule
%\mc{5}{c}{\tsc{Sula-Buru}}\\
\midrule
\ili{Sanana}	&	\cg{(mora)}	&	\cg{(fasina)}	&	{\hr}uya	&	{\hr}yaa (*l > {\0}/a{\gap}a)\\
\ili{Lisela}	&	\cg{(moda)}	&	{\hr}fula\tbr{n}	&	\cg{(dekat)}	&	\cg{(tuhu-n)}\\
\ili{Buru}	&	{\hr} aŋi\tbr{n}	&	{\hr}fula\tbr{n}	&	\cg{(dekat)}	&	\cg{(toho-n)}\\
\ili{Ambelau}	&	{\hr} aŋi	&	{\hr}hula	&	{\hr}ula\tbr{n}	&	{\hr}lah-lea\su{†}\\
%\midrule
%\mc{5}{c}{\tsc{\ili{Ambon}-Seram}}\\
\midrule
\ili{Kayeli}	&	{\hr} ani\tbr{n}	&	{\hr}bula\tbr{n}	&	{\hr}ula\tbr{n}	&	{\hr}lala\tbr{n}\\
\ili{Larike}	&	{\hr} ani\tbr{n}-u	&	{\hr}huda\tbr{n}-u	&	{\hr}uda\tbr{n}-u	&	{\hr}lala\tbr{n}-u\\
\ili{Asilulu}	&	{\hr} ani\tbr{n}	&	{\hr}hula\tbr{n}	&	{\hr}ula, ula\tbr{n}	&	{\hr}lala\tbr{n}\\
\ili{Paulohi}	&	{\hr} ani\tbr{n}-e	&	{\hr}hula\tbr{n}-e	&	\cg{(kia)}	&	{\hr}laa kine\\
\ili{Sepa}	&	{\hr}yani\tbr{n}-o	&	{\hr}hia\tbr{n}-o	&	\cg{(kiya)}	&	{\hr}laa tina\\
N. \ili{Alune}	&	\cg{(sanut-e)}	&	{\hr}bula\tbr{n}-e	&	{\hr}ula\tbr{n}-e	&	{\hr}lala\tbr{n}-e\\
\ili{Yalahatan}	&	\cg{(yolil-e)}	&	{\hr}fula\tbr{n}-e	&	{\hr}ura\tbr{n}-e	&	{\hr}rara\tbr{n}-e\\
S. \ili{Nuaulu}	&	\cg{(ihut-e)}	&	{\hr}huna\tbr{n}-e	&	{\hr}ua\tbr{n}-e	&	\cg{(arena)}\\
\ili{Liana-Seti}	&	\cg{(simala)}	&	\cg{(umala)}	&	\cg{(roa)}	&	{\hr}lala\\ \midrule
\ili{Banda}	&	{\hr} ani\tbr{n}	&	{\hr}ɸula\tbr{n}	&	\cg{(wun)}	&	{\hr}era\tbr{n}\\
%\midrule
%\mc{5}{c}{\tsc{Seram-Tanimbar-Bomberai}}\\
\midrule
\ili{Masiwang}	&	{\hr}yaki\tbr{n}	&	{\hr}hula\tbr{n}	&	\cg{(roti)}	&	{\hr}lolo\tbr{n}\\
\ili{Bobot}	&	{\hr}yaki\tbr{n}	&	{\hr}vula\tbr{n}	&		&	{\hr}lolo\tbr{n}\\
\ili{Geser}-\ili{Gorom}	&	{\hr} aŋi\tbr{n}	&	\hp{*v}ula\tbr{n}	&	{\hr}ura\tbr{n}	&	{\hr}lala\tbr{n}\\
\ili{Kowiai}	&	\cg{(lor)}	&	{\hr}ɸura\tbr{n}	&	\cg{(oma)}	&	{\hr}rara\tbr{n}\\
\ili{Yamdena}	&	\cg{(mnaur)}	&	{\hr}bula\tbr{n}	&	{\hr}uda\tbr{n}	&	{\hr}dal\tbr{n}-ene\\
\ili{Fordata}	&	\cg{(nait)}	&	{\hr}vula\tbr{n}	&	\cg{(daʔut)}	&	\cg{(liŋaʔan)}\\
\ili{Kei}	&	\cg{(niit, yaran)}	&	{\hr}vua\tbr{n} (*l > {\0})	&	\cg{(doʔot)}	&	\cg{(ded)}\\
\ili{Irarutu}	&	\cg{(nóɸə)}	&	\cg{(sibá)}	&	\cg{(syem)}	&	\cg{(rarúmə)}\\
		\lspbottomrule
	\end{tabular}
		\begin{tablenotes}\tnsize
			\item[†] {The first part of \ili{Ambelau} \it{lah} is probably a reflex of *zalan,
								since *z > \it{l}, and *l > \it{l, h}.}
		\end{tablenotes}
	\end{threeparttable}
\end{table}

\trf{07-02} shows retention of final *t
and *R in geographical central Maluku languages across several subgroups.
Retentions of final *k,
*l, *m, *p, *s,
*j, *q and possibly *h are also found,
but to a lesser degree.
Note that final *R is lost in \tsc{Sula-Buru},
but retained in \tsc{\ili{Ambon}-Seram}
and \tsc{Seram-Tanimbar-Bomberai}.

\begin{table}[h]\tcp{Geographical central Maluku retention of *t and *R /{\gap\#}}{07-02}
%\begin{adjustbox}{width=\textwidth}
	\st{3.5pt}\begin{tabular}{llllll}\lsptoprule
local gl.	&	sky	&	west monsoon	&	\ili{east} monsoon	&	sail	&	coconut\\
PMP	&	*laŋi\tbr{t}	&	*habaRa\tbr{t}	&	*timu\tbr{R}	&	*laya\tbr{R}	&	*ñiu\tbr{R}/**niwə\tbr{R}\\
%\midrule
%\mc{6}{c}{\tsc{Sula-Buru}}\\
\midrule
\ili{Sanana}	&	{\hr}lan	&	\cg{(bara)} <Mly.	&	{\hr}timu	&	\cg{(soba)}	&	{\hr}nui \it{(met.)}\\
\ili{Buru}	&	{\hr}laŋi\tbr{t}	&	{\hr}ful faha\tbr{t}	&	{\hr}ful timo	&	{\hr}laa	&	{\hr}niwe\\
\ili{Ambelau}	&	{\hr}lani\tbr{r}-e	&	{\hr}baha\tbr{r}-e	&	{\hr}rimu	&	{\hr}haa	&	{\hr}niβe\\
%\midrule
%\mc{6}{c}{\tsc{\ili{Ambon}-Seram}}\\
\midrule
\ili{Kayeli}	&	{\hr}lani\tbr{t}	&		&		&	{\hr}laa-ni	&	{\hr}niwe\tbr{l}-e\\
\ili{Larike}	&	{\hr}lani\tbr{t}-a	&	{\hr}hala\tbr{t}-a	&	{\hr}timu\tbr{d}-u	&	{\hr}lae\tbr{d}-u	&	{\hr}niwe\tbr{d}-u\\
\ili{Asilulu}	&	{\hr}lani\tbr{t}	&	{\hr}hala\tbr{t}	&	{\hr}timu\tbr{l}	&	{\hr}laa\tbr{l}	&	{\hr}niwe\tbr{l}\\
\ili{Paulohi}	&	{\hr}lani\tbr{t}-e	&	{\hr}hala\tbr{t}-e	&	{\hr}timu\tbr{l}-e	&	{\hr}laa\tbr{l}-e	&	\\
\ili{Sepa}	&	{\hr}laan\tbr{t}-o	&	{\hr}hara\tbr{t}-o	&	{\hr}timi-o	&	{\hr}lia\tbr{l}-o	&	{\hr}nuwe\tbr{l}-o\\
N. \ili{Alune}	&	{\hr}lani\tbr{t}-e	&	{\hr}bala\tbr{t}-e	&	{\hr}timu\tbr{l}-e	&	{\hr}lae\tbr{l}-e	&	{\hr}nikwe\tbr{l}-e\\
\ili{Yalahatan}	&	{\hr}lan\tbr{t}-e	&	{\hr}hala\tbr{t}-e	&	{\hr}timi\tbr{l}-e	&	{\hr}leʔa\tbr{l}-e	&	{\hr}nie\tbr{l}-e\\
S. \ili{Nuaulu}	&	{\hr}nan\tbr{t}-e	&	{\hr}hana\tbr{t}-e	&	\cg{(ihat-e)}	&	{\hr}na\tbr{n}-e	&	{\hr}nio\tbr{n}-e\\
\ili{Liana-Seti}	&	\cg{(lea)}	&	{\hr}fara\tbr{t}-a	&	{\hr}tim\tbr{l}-a	&	{\hr}laiya\tbr{l}-a	&	{\hr}nua \\ \midrule
\ili{Banda}	&	{\hr}laŋi\tbr{t}	&	{\hr}fara\tbr{t}	&	{\hr}timu\tbr{r}	&	{\hr}lea\tbr{r}	&	{\hr}nue\tbr{r}\\
%\midrule
%\mc{6}{c}{\tsc{Seram-Tanimbar-Bomberai}}\\
\midrule
\ili{Masiwang}	&	{\hr}laki\tbr{t}	&	{\hr}wara\tbr{t}	&	{\hr}timi &	{\hr}lai	&	\\
\ili{Bobot}	&	{\hr}laki\tbr{t}	&		&		&		&	\\
\ili{Geser}	&	{\hr}laŋi\tbr{t}	&	{\hr}wara\tbr{t}	&	{\hr}timo\tbr{r}	&	{\hr}la\tbr{r}	&	{\hr}niu, niu\tbr{r}\\
\ili{Kowiai}	&	{\hr}raŋgi\tbr{t}	&	\cg{(mondar)}	&	{\hr}timo\tbr{r}	&	\cg{(ambaʔ)}	&	{\hr}niu, niu\tbr{r}\\
\ili{Yamdena}	&	{\hr}laŋi\tbr{t}	&	{\hr}bara\tbr{t}	&	{\hr}timu\tbr{r}	&	{\hr}la\tbr{r}e	&	{\hr}nu\tbr{r}\\
\ili{Fordata}	&	{\hr}lani\tbr{t}	&	{\hr}vara\tbr{t}	&	{\hr}timu\tbr{r}	&	{\hr}laa\tbr{r}	&	{\hr}nuu\tbr{r}\\
\ili{Kei}	&	{\hr}lani\tbr{t}	&	{\hr}vara\tbr{t}	&	{\hr}timu\tbr{r}	&	{\hr}laa\tbr{r}	&	{\hr}nuu\tbr{r}\\
\ili{Irarutu}	&	{\hr}ragə\tbr{t}ə	&		&	{\hr}timú\tbr{r}ə	&	\cg{(marú)}	&	\cg{(umagí)}\\
		\lspbottomrule
	\end{tabular}
%\end{adjustbox}
\end{table}

There is also some confusion in both descriptive
and comparative works, in that a number of so-called “retentions”
of final *C{\#} (such as *j
and *t) are actually productive suffixes in the synchronic morphosyntax
of the languages of geographical central Maluku.\footnote{\label{fn:Collins-Final-t}
		For example, \citet[32]{Collins1981} argued that PMP *t split in \ili{Taliabu} becoming
		palatal \it{c} [tsʸ] in final position.
		However, it is actually a productive nominal suffix \it{-c}.
		(\tsc{Sula-Buru} and \tsc{\ili{Ambon}-Seram} cognates for the same
		words all have the equivalent nominal \it{-t(e)} suffix,
		and roots that take it can often occur with other suffixes as
		well.) \citet[34]{Collins1981} also says “Based on the Taliabo
		and \ili{Buru} entries for *waRej it is argued that in BST [\ili{Buru}-Sula-Taliabo]
		*j in final position became *t {\ldots}[but] were inexplicably
		lost” in other words that end in *j.
		He cites *waRəj ‘vine,
		rope’ > \ili{Buru} \it{wahet}
		and \ili{Taliabu} \it{tuka wahoc} ‘guts’ as evidence.
		But \ili{Buru} has \it{wahe-t}
		and \it{wahe-n}, both of which are entirely consistent with those
		productive suffixes in \ili{Buru} (see \citealp[987]{GrimesGrimes2020})
		and \ili{Taliabu} \it{waho} ‘rattan’ is cited as showing “loss of final
		\it{c}” (p.43, footnote 24).
		\citet[36]{Collins1981} cites “the merger of *j
		and *t in final position” as one of the innovations that allows
		including \ili{Ambelau} with his BST [\ili{Buru}-Sula-Taliabo] under West \ili{Central} Maluku.
		In his summary (\citeyear[38]{Collins1981}) he says “In final position *j became
		*t in PWCM [Proto-West \ili{Central} Maluku].”
		\citet[39--40]{Collins1981} cites similar reflexes from *waRəj from
		East \ili{Central} Maluku languages to support the claim that “at l\ili{east}
		in one entry PECM [Proto-East \ili{Central} Maluku] shifted final *j to *t.”
		In \citet[40]{Collins1981} he lists *j > \it{t}/{\gap}{\#}
		as an innovation defining Proto-\ili{Central} Maluku.
		However, many \tsc{\ili{Ambon}-Seram} languages also have the
		productive \it{-t(e)} nominal suffix.
		In footnote 8 (\citeyear[42]{Collins1981}) Collins does provide a caveat to his
		claims: “There is a possibility that this merger to \it{t} in
		final position is related to the grammatical suffix \it{-t} used
		to mark dependent nouns or verbs,” but then dismisses it.
		It is a productive suffix across the region and beyond.
		See \citet[30--33]{Collins1983} for several languages; \citet[xxx]{Collins2003a}
		for \ili{Asilulu} nominal \it{-t}; \citet{Grimes1991c} for \ili{Buru} nominal
		\it{-t}; \citet{Bolton1990} for South \ili{Nuaulu} nominal \it{-te;}
		\citet{Travis2019} for \ili{Kei} \it{-t}.
		Beyond geographical central Maluku see \citet{Tamelan2021} for
		\ili{Dela} nominal \it{-t/-ʔ/-k}.}
There may be phonotactic pressure in some languages to lose historical
root-final consonants to make way for pronominal enclitics,
possessive enclitics, nominalisers, noun markers, detransitivisers,
and applicatives, which often take the shape of \it{-n,
-ɲ, -ni, -i, -e,
-t, -te, -c}, and \it{-k} in many languages of geographical central Maluku.
Some of these are illustrated in chapters \ref{ch:Tal}--\ref{ch:Ban},
and discussed in the relevant grammars in the region cited in the Bibliography.
A similar pattern was noted by \citet{Jonker1906} for languages
in the Timor region.
However the final consonants illustrated in \trf{07-01}
and \trf{07-02} are not readily explained by the grammatical functions
of those suffixes and are best analysed as historical retentions.
Investigators are easily confused about which are historical
retentions, and which are productive suffixes.

Looking at the wider Austronesian picture,
we note that the nearby \tsc{\ili{Bungku}-Tolaki} languages of south\ili{east}
Sulawesi -- west of the Blust line -- have lost many final consonants
across the board for certain words in all languages,
including for *aŋin, *quzan,
*bulan, *laŋit, *habaRat, *timuR \citep[103--178]{Mead1999}.\footnote{
		\citet{Mead1996}
		shows that final consonants need to be reconstructed for Proto-\ili{Bungku}-\ili{Tolaki},
		even if they are not found on specific lexical items.}
Yet these final consonants missing on certain words in \tsc{\ili{Bungku}-Tolaki}
languages are actually retained in many languages of geographical central Maluku.
So a generalisation that the languages \ili{east} of the Blust line
are losing or have lost historical final consonants,
with the implication that languages west of the line have not,
is not well grounded in the data.


